\section{L6 -- Memory}

\subsection{Cache \& Paging}

\mult{2}

\begin{definition}{Cache}\\
    Fast memory near/in the CPU, \important{transparent} to the programmer.
    \begin{itemize}
        \item hit $\to$ served fast
        \item miss $\to$ fetch a whole \emph{line} from memory (also cached)
    \end{itemize}
\end{definition}

\begin{definition}{Paging}\\
    Address space $\to$ pages (4K/2M); physical $\to$ frames.
    \begin{itemize}
        \item load only used pages (any frame)
        \item \textbf{resident set} = pages in memory
        \item \textbf{working set} = pages in use now
    \end{itemize}
\end{definition}

\multend

\begin{concept}{Why a Hierarchy}\\
    The CPU is far faster than RAM (the bottleneck), so caches + paging hide latency and use main memory efficiently.
\end{concept}

\begin{example2}{Ex.\ S4/5 -- Explain cache \& paging}\\
    What is a cache (hit/miss)? What is paging (resident vs working set)?
    \tcblower
    Cache = fast transparent memory; miss $\to$ fetch a line. Paging = pages$\to$frames, only used pages resident; \important{resident} = in memory, \important{working} = in use now.
\end{example2}

\subsection{The Memory Access Walk}

\begin{definition}{Access Chain}\\
    CPU pipeline (IF$\cdot$ID$\cdot$EX$\cdot$MEM$\cdot$WB) $\to$ Cache $\to$ \textbf{MPU} (rights) $\to$ \textbf{MMU} (\textbf{TLB} = recent mappings; else page-table walk) $\to$ \textbf{page fault} $\to$ MM software.
\end{definition}

\begin{KR}{Recipe: Trace a Memory Access}\\
    \textbf{Pipeline} stage requests an address.
    \textbf{Cache} hit $\to$ done; miss $\to$ continue.
    \textbf{MPU} access allowed?
    \textbf{MMU/TLB} hit $\to$ frame; miss $\to$ walk page tables.
    \textbf{Page fault} empty $\to$ MM loads frame from disk, updates page dir.
    \textbf{Fill} TLB caches mapping; cache reads a line; write sets \important{dirty bit}.
\end{KR}

\begin{example2}{Ex.\ S6--12 -- Worked code trace}\\
    Walk \texttt{mov ea,var1; mov eb,var2; add ea,eb; mov var3,ea}. Why first fetch slow, next fast?
    \tcblower
    First fetch (\texttt{0x100}): cache miss $\to$ MPU $\to$ MMU $\to$ TLB miss $\to$ walk empty $\to$ \important{page fault} (load frame \#2). Fill = whole line (I1--I4) + TLB mapping $\Rightarrow$ \texttt{0x104},\texttt{0x108} are \important{cache hits}. \texttt{var1} (\texttt{0x400}): 2nd page fault (frame \#8); \texttt{var2} = hit; \texttt{mov var3} = write $\to$ \important{dirty bit}.
    % INSERT IMAGE: L6 slide 6/7 -- Simple Computer System walk
\end{example2}

\subsection{Thrashing \& Context-Switch Cost}

\mult{2}

\begin{definition}{Thrashing}\\
    \begin{itemize}
        \item \textbf{cache}: many lines map to one slot $\to$ evict/miss loop
        \item \textbf{disk}: too few frames $\to$ page evict/reload loop
    \end{itemize}
\end{definition}

\begin{concept}{Switch Cost}\\
    \begin{itemize}
        \item cache \important{high} (flush + write-back)
        \item TLB \important{high} (flush)
        \item MPU/page-tables medium; pipeline low
    \end{itemize}
\end{concept}

\multend

\begin{example2}{Ex.\ S13--15 -- Thrashing \& switch cost}\\
    Explain both thrashing types and the cost of cache-oblivious code.
    \tcblower
    Cache thrash = repeated misses (same slot); disk thrash = repeated page evict/reload. Switch cost dominated by cache+TLB flush; CPI spikes then re-warms.
\end{example2}

\subsection{Cache Writeback \& Multicore}

\begin{definition}{Writeback (dirty bit)}\\
    A dirty line is written back when (1) the controller decides, or (2) it needs the line (eviction). Then dirty$\to$0; eviction sets valid$\to$0, refill, valid$\to$1.
\end{definition}

\begin{KR}{Recipe: Will a Writeback Happen?}\\
    Check the dirty bit: dirty + (controller decision or eviction) $\to$ write back first. Clean line $\to$ just refill.
\end{KR}

\begin{example2}{Ex.\ S16--18 -- Writeback \& multicore L1/L2}\\
    When does a dirty line write back? L1 read/write-back on shared L2?
    \tcblower
    See recipe. Multicore: read pulls L2$\to$L1; \important{flushing L1 requires flushing L2}; L1$\to$L2 write-back may trigger L2$\to$memory; coherence keeps other L1s consistent.
\end{example2}

\subsection{Cache Coherence}

\mult{2}

\begin{definition}{Coherence}\\
    Visibility of a write across cores. A write \emph{will} become visible (order to one location preserved) -- says nothing about \emph{when}; not in ISA; necessary not sufficient.
\end{definition}

\begin{definition}{Consistency}\\
    Ordering of accesses to \emph{different} locations as seen by another core.
    \begin{itemize}
        \item sequential / processor / release
        \item release: order only at barriers
    \end{itemize}
\end{definition}

\multend

\mult{2}

\begin{concept}{Snooping}\\
    Bus broadcast; controllers listen, invalidate/update. Non-scalable.
\end{concept}

\begin{concept}{Directory}\\
    Point-to-point; owner invalidates sharers, waits ack, writes, becomes owner. Scalable.
\end{concept}

\multend

\begin{definition}{False Sharing}\\
    Two cores write \emph{different} vars in the \emph{same line} $\to$ line ping-pongs (useless coherence traffic). Fix: pad/align to separate lines.
\end{definition}

\begin{KR}{Recipe: Diagnose Coherence / False Sharing}\\
    Same line in $\geq$2 caches + a write $\to$ invalidation traffic. Same var = true sharing; different vars = \important{false sharing} $\to$ pad/align.
\end{KR}

\begin{example2}{Ex.\ S19/20 -- Coherence vs consistency}\\
    Define both; what does 0-1-1-2 / 0-2-1-2 show?
    \tcblower
    Two cores observe writes in \emph{different} order $\to$ a protocol must make P2 see P1's order (coherence). Consistency = cross-location ordering (sequential/processor/release).
\end{example2}

\begin{example2}{Ex.\ S21--23 -- Protocols \& false sharing}\\
    Outline snooping vs directory; explain false sharing.
    \tcblower
    Snooping = bus broadcast invalidation (e.g.\ P1 writes \texttt{var1} $\to$ P2 invalidated), non-scalable. Directory = owner-based, scalable. False sharing = different vars, same line $\to$ ping-pong; fix by padding.
\end{example2}

\begin{example2}{Ex.\ S24 -- ASID / CPU pinning}\\
    How is cache/TLB state preserved across switches?
    \tcblower
    Physical-address tagging or an \important{ASID} lets one process's lines survive while another runs $\Rightarrow$ cheap to re-schedule on the same core (\important{CPU pinning / affinity}).
\end{example2}

\subsection{Memory Architectures}

\begin{definition}{UMA / NUMA / SUMA}\\
    \begin{itemize}
        \item \textbf{UMA}: equal access, $\sim$4-CPU limit, L2/L3 buffer
        \item \textbf{NUMA}: fast local, slow remote via interconnect (ccNUMA)
        \item \textbf{SUMA}: interleaving $\to$ best-case UMA
    \end{itemize}
    % INSERT IMAGE: L6 slides 25-28 -- UMA/NUMA/SUMA diagrams
\end{definition}

\begin{KR}{Recipe: Classify a Memory Architecture}\\
    equal access $\to$ UMA; fast-local/slow-remote $\to$ NUMA; interleaved-to-fake-contiguous $\to$ SUMA.
\end{KR}

\begin{example2}{Ex.\ S25--29 -- Architectures \& OS support}\\
    Bus limits? UMA/NUMA/SUMA? How does Linux handle NUMA?
    \tcblower
    Bus latency grows with length $\to$ $\sim$4-CPU limit. UMA/NUMA/SUMA as above. Linux maps HW \emph{cells} $\to$ \emph{nodes}, keeps work+memory local (affinity), \important{zonelists} by NUMA distance; \texttt{taskset}/\texttt{numactl}/\texttt{sched\_setaffinity}.
\end{example2}
